% Generated draft: requires mathematical and source review.
\documentclass[11pt]{article}
\usepackage[margin=25mm]{geometry}
\usepackage{fontspec}
\setmainfont{Latin Modern Roman}
\usepackage{amsmath,amssymb,mathtools,unicode-math}
\setmathfont{Latin Modern Math}
\usepackage{enumitem,needspace,xurl,hyperref}
\setlength{\parindent}{0pt}
\setlength{\parskip}{5pt}
\setlength{\emergencystretch}{3em}
\newcommand{\N}{\mathbb{N}}
\newcommand{\Z}{\mathbb{Z}}
\newcommand{\Q}{\mathbb{Q}}
\newcommand{\R}{\mathbb{R}}
\newcommand{\C}{\mathbb{C}}
\newcommand{\F}{\mathbb{F}}
\newcommand{\A}{\mathbb{A}}
\newcommand{\PP}{\mathbb{P}}
\newcommand{\E}{\mathbb{E}}
\newcommand{\eps}{\varepsilon}
\newcommand{\dd}{\,\mathrm{d}}
\begin{document}
\section*{UnsolvedMath ID: 30005945}
\textbf{Source problem code:} catalogue entry
\subsection{Definitions and mathematical statement}
Let \(Y\) be a smooth complex projective Fano variety, meaning that its anticanonical line bundle \(\mathcal O_Y(-K_Y)\) is ample. Here \(K_Y\) denotes a canonical divisor, and \(|-K_Y|\) denotes the complete linear system of effective divisors linearly equivalent to \(-K_Y\). A smooth anticanonical divisor is a divisor \(D\in|-K_Y|\) that is smooth over \(\C\).

The supplied characterization problem is to determine exactly which such varieties \(Y\) satisfy
\[
\exists\,D\in|-K_Y|\text{ smooth},\qquad
\exists\,f:(\mathcal Y,\mathcal D)\longrightarrow(S,0)
\]
where \(f\) is a maximal log-smooth degeneration with generic fiber \((Y,D)\), in the sense intended by the cited source. The quantification is existential over both the divisor and the degeneration; no particular divisor is fixed in advance.

\emph{Unresolved qualifications.} The supplied record does not define ``maximal log-smooth degeneration,'' specify the logarithmic structures or the pointed base \((S,0)\), or explain the identification of its generic fiber with the complex pair \((Y,D)\). Consequently, this is a source-qualified characterization target: an intrinsic membership criterion cannot be formulated unambiguously from this record alone. No additional conditions on the base, total space, or special fiber are imposed in this draft.
\subsection{Short English statement}
Characterize the smooth complex projective Fano varieties that admit at least one smooth anticanonical divisor and a maximal log-smooth degeneration of the resulting pair. The source-specific meaning of the degeneration, especially ``maximal,'' must be recovered before the formulation is mathematically complete.
\subsection{Sources}
\begin{itemize}
\item Local UnsolvedMath registry, supplied source record: permanent ID \(30005945\), source problem code \texttt{catalogue entry}, title ``Fano Varieties with Maximal Log-Smooth Degenerations.''
\item Supplied report citation: \emph{Real and Logarithmic Enumerative Geometry} (2024), Oberwolfach Reports, pp.~1285--1338, DOI \texttt{10.4171/owr/2024/23}. Supplied source URL and literature-sources link, labeled ``Oberwolfach Report 23/2024'': \url{https://doi.org/10.4171/owr/2024/23}.
\item The original excerpt asks how restrictive the existence of \((Y,D,f)\) is for given \(Y\), and associates the discussion with ``Enumerative intrinsic mirror symmetry,'' Michael van Garrel. The registry does not separately identify a proposer.
\item Literature assessment dated exactly \texttt{2026-08-22}: the supplied curation labels the problem open, reports that the source explicitly asks about the restrictiveness of a maximal log-smooth degeneration, and reports that no resolution was located. This is supplied dated curation, not an independently verified literature assessment. The background requests independent verification of time-sensitive status and cited claims.
\item The registry labels the statement \texttt{corrected\_verified} and reports that it was checked against the official Oberwolfach report for preservation of hypotheses, notation, and target. That verification claim is supplied provenance and has not been independently checked for this draft.
\end{itemize}
\end{document}
